% !TEX program = pdflatex
\documentclass[10pt]{article}
\usepackage[utf8]{inputenc}\usepackage[T1]{fontenc}\usepackage[english]{babel}
\usepackage{lmodern,microtype}\usepackage[margin=.72in]{geometry}
\usepackage{amsmath,amssymb,amsthm,mathtools,booktabs,array,xcolor,enumitem}
\usepackage[most]{tcolorbox}\usepackage{fancyhdr}\usepackage[hidelinks]{hyperref}
\definecolor{CourseBlue}{HTML}{1A3C68}\definecolor{CourseGreen}{HTML}{216E39}
\definecolor{CourseGold}{HTML}{B8860B}\definecolor{CourseRed}{HTML}{9A2C2C}
\definecolor{SoftBlue}{HTML}{F2F6FA}\definecolor{SoftGreen}{HTML}{F1F8F3}
\definecolor{SoftGold}{HTML}{FFF8E8}\definecolor{SoftRed}{HTML}{FFF3F3}
\setlength{\parindent}{0pt}\setlength{\parskip}{.32em}
\setlist[itemize]{leftmargin=1.5em,itemsep=.12em,topsep=.2em}
\setlist[enumerate]{leftmargin=1.7em,itemsep=.12em,topsep=.2em}
\pagestyle{fancy}\fancyhf{}\lhead{Stochastic Calculus and Option Pricing}
\rhead{Lecture 11 Notes}\cfoot{\thepage}\renewcommand{\headrulewidth}{.3pt}
\newcommand{\E}{\mathbb E}\newcommand{\Q}{\mathbb Q}\newcommand{\cF}{\mathcal F}
\newcommand{\dd}{\,\mathrm d}\newcommand{\esssup}{\operatorname*{ess\,sup}}
\theoremstyle{definition}\newtheorem{definition}{Definition}[section]
\newtheorem{proposition}[definition]{Proposition}
\newtcolorbox{boardexample}[1]{enhanced,breakable,colback=SoftBlue,colframe=CourseBlue,boxrule=.55pt,arc=1mm,title={Example --- #1},fonttitle=\bfseries}
\newtcolorbox{studentexercise}[1]{enhanced,breakable,colback=SoftGold,colframe=CourseGold,boxrule=.55pt,arc=1mm,title={Exercise --- #1},fonttitle=\bfseries}
\newtcolorbox{shortsolution}[1][]{enhanced,breakable,colback=SoftGreen,colframe=CourseGreen,boxrule=.55pt,arc=1mm,title={Short solution#1},fonttitle=\bfseries}
\newtcolorbox{keyidea}[1]{enhanced,breakable,colback=white,colframe=CourseGreen,boxrule=.7pt,arc=1mm,title={#1},fonttitle=\bfseries}
\newtcolorbox{warningbox}[1]{enhanced,breakable,colback=SoftRed,colframe=CourseRed,boxrule=.65pt,arc=1mm,title={#1},fonttitle=\bfseries}
\title{\textbf{Change of Numeraire}\\\large Lecture 11 Notes --- Units, Measures, Forward Pricing, and FX}
\author{}\date{}
\begin{document}\maketitle\vspace{-1.4em}
\begin{keyidea}{Goals and logical structure}
Risk-neutral valuation already uses a unit: the money-market account. We
change that unit deliberately:
\[
\text{positive traded numeraire}\longrightarrow\text{matching measure}
\longrightarrow\text{martingale price ratio}\longrightarrow\text{simpler pricing}.
\]
The price does not change; only its representation does.
\end{keyidea}

\section{Numeraires, measures, and ratios}
\subsection{The general pricing formula}
With money-market account \(B_t\), a payoff \(X_T\) is worth
\[
V_t=B_t\E^\Q\!\left[\frac{X_T}{B_T}\middle|\cF_t\right].
\]
\begin{definition}[Numeraire]
A numeraire \(N_t\) is a strictly positive traded asset used as the unit of
account. The \(N\)-denominated price of an asset \(S\) is \(S_t/N_t\).
\end{definition}
\begin{proposition}[Change-of-numeraire density]
Starting under \(\Q\), the measure associated with \(N\) is
\[
\left.\frac{\dd\Q^N}{\dd\Q}\right|_{\cF_t}
=\frac{N_t/B_t}{N_0/B_0}.
\]
Under \(\Q^N\), every traded price divided by \(N_t\) is a martingale.
\end{proposition}
\begin{keyidea}{One reusable formula}
\[
\boxed{V_t=N_t\E^{\Q^N}\!\left[\frac{X_T}{N_T}\middle|\cF_t\right].}
\]
Divide the payoff by the numeraire, take expectation under its associated
measure, then restore the current unit.
\end{keyidea}
\begin{warningbox}{Unit and measure must move together}
Dividing by a new numeraire while retaining the old measure generally leaves
a drift. The measure change reweights scenarios but does not create or
remove possible market scenarios.
\end{warningbox}
\begin{studentexercise}{a bond numeraire}
If \(N_t=P(t,T)\), what quantity is a martingale and what is the measure
called?
\end{studentexercise}
\begin{shortsolution}
For every traded asset \(S\), \(S_t/P(t,T)\) is a martingale under the
\(T\)-forward measure \(\Q^T\).
\end{shortsolution}

\subsection{Ratios need It\^o corrections}
If
\[
\frac{\dd S_t}{S_t}=\mu_S\dd t+\sigma_S\dd W_t,\qquad
\frac{\dd N_t}{N_t}=\mu_N\dd t+\sigma_N\dd W_t,
\]
then
\[
\frac{\dd(S_t/N_t)}{S_t/N_t}
=\bigl(\mu_S-\mu_N+\sigma_N^2-\sigma_S\sigma_N\bigr)\dd t
+(\sigma_S-\sigma_N)\dd W_t.
\]
The drift is not just \(\mu_S-\mu_N\): quadratic variation matters.

\section{Forward measures match payment dates}
\subsection{A bond absorbs discounting}
Choose the bond \(P(t,T)\). Since \(P(T,T)=1\),
\[
\boxed{V_t=P(t,T)\E^{\Q^T}[X_T\mid\cF_t].}
\]
For a non-dividend-paying asset,
\[
F(t,T)=\frac{S_t}{P(t,T)}
\]
is a \(\Q^T\)-martingale.
\begin{boardexample}{a zero-value forward}
A forward delivering one share for \(K\) has value
\[
V_t=S_t-KP(t,T)=P(t,T)(F(t,T)-K).
\]
If \(S_0=100\) and \(P(0,T)=.95\), the fair delivery price is
\(F(0,T)=105.26\).
\end{boardexample}
\subsection{A factor-cancelling choice}
If \(X_T=S_TY_T\), the stock numeraire yields
\[
V_t=S_t\E^{\Q^S}[Y_T\mid\cF_t].
\]
The candidate must be strictly positive; an arbitrary payoff is not a
numeraire.
\begin{studentexercise}{first candidate}
A claim pays at \(T\) and depends on future bond prices. Which numeraire is
the natural first candidate?
\end{studentexercise}
\begin{shortsolution}
\(P(t,T)\): it matches the payment date and removes discounting from the
expectation under \(\Q^T\).
\end{shortsolution}

\section{FX: two currencies, two units}
Let \(X_t\) be domestic currency per foreign currency unit and
\[
B_d(t)=e^{r_dt},\qquad B_f(t)=e^{r_ft}.
\]
The domestic value of the foreign bank account is \(X_tB_f(t)\). Therefore,
under the domestic measure,
\[
\boxed{\frac{\dd X_t}{X_t}=(r_d-r_f)\dd t+\sigma\dd W_t^{\Q^d}.}
\]
Foreign interest plays the role of a continuous dividend yield.
\begin{boardexample}{FX forward}
\[
F_{0,T}=X_0e^{(r_d-r_f)T}.
\]
For \(X_0=1.10\) USD/EUR, \(r_d=5\%\), \(r_f=3\%\), and \(T=1\),
\(F_{0,1}=1.1222\) USD/EUR.
\end{boardexample}

\subsection{FX options}
A domestic call on one foreign-currency unit pays \((X_T-K)^+\):
\[
C_t=X_te^{-r_f\tau}\Phi(d_1)-Ke^{-r_d\tau}\Phi(d_2),
\]
\[
d_1=\frac{\log(X_t/K)+(r_d-r_f+\frac12\sigma^2)\tau}{\sigma\sqrt{\tau}},
\qquad d_2=d_1-\sigma\sqrt{\tau}.
\]
Parity makes the two discount factors visible:
\[
\boxed{C_t-P_t=X_te^{-r_f\tau}-Ke^{-r_d\tau}.}
\]
\begin{warningbox}{Keep currencies explicit}
The spot term represents foreign cash and is discounted at \(r_f\). The
strike is domestic cash and is discounted at \(r_d\).
\end{warningbox}

\section{Synthesis and practice}
\begin{keyidea}{Checklist}
\begin{enumerate}
\item State payoff date and currency.
\item Select a strictly positive traded numeraire.
\item Use its associated measure.
\item Price the normalized payoff and restore the current unit.
\item Check against a forward relation or money-market pricing.
\end{enumerate}
\end{keyidea}
\begin{studentexercise}{final check}
Why is it false that a change of numeraire changes the economic price?
\end{studentexercise}
\begin{shortsolution}
The claim has one no-arbitrage value. The unit, equivalent probability
weights, and martingale ratio change; the economic price does not.
\end{shortsolution}
\end{document}
